% !PRINTED 603
\weberpage{603}
%<<CONT

Wir nehmen nach \S\,161, 4. in der Klasse $(A_i)^{-1}$ ein zu $M$ teilerfremdes Ideal $\mathfrak{m}$ und lassen in

\weq{5}{\mathfrak{m}\,\mathfrak{a}_i = \overline{\alpha}}

$\mathfrak{a}_i$ die Ideale der Klasse $\overline{A}_i$, und folglich $\overline{\alpha}$ die durch $\mathfrak{m}$ teilbaren Ideale der Hauptklasse durchlaufen; ist dann

\[
N(\mathfrak{a}_i) \leqq \nu,
\]

so ist

\[
N(\mathfrak{m})\,N(\mathfrak{a}_i) = N(\alpha) \leqq \nu\,N(\mathfrak{m}).
\]

Bedeutet wie früher $T$ die Anzahl der durch $\mathfrak{m}$ teilbaren Ideale der Hauptklasse, deren Norm nicht größer als $t$ ist, so ist, wenn $t = \nu N(\mathfrak{m})$ gesetzt wird,

\[
Z(\nu) = T,
\]

und aus der Formel \S\,162, (6) ergibt sich

\weq{6}{Z(\nu) = g\nu + R_\nu\sqrt{\nu},}

worin $g$ eine Zahl, die für alle Klassen $\overline{A}_i$ dieselbe ist, und $R_\nu$ eine Funktion von $\nu$ bedeutet, die mit unendlich wachsendem $\nu$ nicht unendlich wird.

Nach der Bedeutung von $a_\nu$ ist

\weq{7}{a_\nu = Z(\nu) - Z(\nu - 1), \quad Z(0) = 0}

und folglich

\weq{8}{S_i = \sum_{0,\infty}^{\nu}\frac{Z(\nu) - Z(\nu - 1)}{\nu^s}.}

Setzen wir für $Z(\nu)$ den Ausdruck (6), so folgt:

\weq{9}{S_i = g\sum\frac{1}{\nu^s} + \sum R_\nu\sqrt{\nu}\left(\frac{1}{\nu^s} - \frac{1}{(\nu+1)^s}\right).}

Das zweite Glied dieser Summe:

\[
\sum R_\nu\sqrt{\nu}\left(\frac{1}{\nu^s} - \frac{1}{(\nu+1)^s}\right)
\]

ist unbedingt konvergent, solange $s > \frac{1}{2}$ ist, und ist daher eine stetige Funktion von $s$, die für $s = 1$ in

\[
\sum_{1,\infty}\frac{R_\nu}{\sqrt{\nu}\,(\nu+1)}
\]

übergeht; denn es ist

\[
\operatorname{Lim}_{\nu=\infty} \nu^{s+\frac{1}{2}} R_\nu \sqrt{\nu}\left(\frac{1}{\nu^s} - \frac{1}{(\nu+1)^s}\right) = s\operatorname{Lim} R_\nu
\]
%CONT>>
